\documentclass[12pt]{article}

% Anything after the % is a comment and is ignored.

% The odd side margin is 1 inch less than the margin on odd numbered
% pages.  This is sufficient for articles, but for books, one often
% sets the odd and even side margins separately.
 
\oddsidemargin0in
\textwidth6.5in

\textheight9in

\usepackage{amsmath,graphicx}

\title{\LaTeX \ essentials}
\author{L. F. Rossi}
\date{\today}

\begin{document}

\maketitle

\section{Introduction}
\label{sec.intro}

\LaTeX \ is a mark-up language for writing mathematics.  Your
priorities should be
\begin{enumerate}

\item Write about your research.\footnote{This really has nothing to
    do with \LaTeX.  You should always be writing because it helps you
    think.}

\item Write your mathematical ideas clearly.  Use your writing to pose
  questions and find pathways to solutions.

\item Continually read and revise your work to make it clearer.

\end{enumerate}

\section{Getting going in \LaTeX}

\LaTeX \ (and \TeX \ from which is came) is
old but very robust.  It's behavior may seem odd at times to users who
are used to modern computers and modern operating systems.  \LaTeX \ was
designed to be useful on machines that are
considerably less capable than what is now available.

\LaTeX \ is free, but there also some commercial packages that provide
nice front-ends for LaTeX.  If you are a linux user, your distribution
probably already includes \LaTeX \ and you are ready to go.  If you are
a Mac user, there is a package called TeXShop that is freely available
to you.  If you are windows user, there are two packages, one called
MikTeX and another called TeXLive, that are available to you.

After you install \LaTeX \ on your computer, you have to find a way to
do things with it.  You will need a plain text editor to create and
edit your tex files.  (I recommend emacs because it is free, powerful
and has very nice features for handling \LaTeX \ files.)  Let's say
your file is called {\tt myfile.tex}.  Typically, this means typing
{\tt latex myfile} from a command line while in the subdirectory
containing your file.  You do not need to include the ``.tex''
extension.  (For instance, on Windows machines, you will need to run
Accessories $>$ Command Prompt, and then ``cd'' to the right
directory.)

If all goes well, you will have created a number of files with
different suffixes like {\tt myfile.aux}, {\tt myfile.log} and {\tt
  myfile.dvi}.  The latter one is very important because {\em dvi}
stands for {\em d}e{\em v}ice {\em i}ndependent.  This file can be
interpreted for any number of devices such as a screen previewer, an
inkjet printer, a laser printer or a jumbo poster printer.  The
important thing is that characters and fonts will look beautiful no
matter what you do with it.  However, you need some tools for viewing
your dvi file, and the good news is that all distributions come with
dvi viewers and assorted tools.  You just have to probe around a bit.

If all does not go well and \LaTeX \ encounters an error.  It will
give you the line number where it encountered problems which is a good
clue for where the error is, and ask you what you want to do.  There
are two useful responses.  The first is ``x'' which is to exit and do
nothing more.  It will not generate any ``dvi'' files for you.  The
second is ``q'' which tells \LaTeX \ to produce a ``dvi'' as far as it
can until it encounters the error.

When it comes time to print your work, your package will come
with some tools for printer dvi files.  A nice way to share documents
is to create files in Adobe's portable document format (pdf).  There
are a number of ways to do this, but the most direct is to use the
command {\tt pdflatex myfile}.

\section{Writing mathematics}

Writing in \LaTeX \ is easy as you can see in Section \ref{sec.intro}.
Writing mathematics is a little more complicated.  One way to write
mathematics is to embed your math in your text as in $f(x) =
\sqrt{1+x+x^2} \neq \sqrt{1} + \sqrt{x} + \sqrt{x^2}$, 
or $\sum_{n=1}^\infty \frac{1}{n^2} =
\frac{\pi^2}{6}$.  Be careful about things like sine and cosine which
really should be set as function: $\sin(x) \neq sin(x)$.  Notice also
that the summation $\sum_{n=1}^\infty$ is different from the capital
Greek symbol $\Sigma_{n=1}^\infty$.  Often, equations should be set by
themselves either numbered
\begin{equation}
f(x) = \int_0^x g(t) dt,
\label{my_equation}
\end{equation}
or without a number
\[
f'(x) = g(x) dt.
\]
(Always, punctuate your equations.)
Referring to equation numbers should be something handled by \LaTeX.
Notice that $\int$ and $\sum$ are special symbols because the
subscripts and superscripts and set in special locations (see
Equation \ref{my_equation}).  Consider this:
\begin{equation}
\sum_{n=1}^\infty \frac{1}{n^2} =
\frac{\pi^2}{6}.
\end{equation}
Notice the difference between the subscripts and superscripts when the
equation is set by itself and when it is embedded in the text.

\section{Learn it just-in-time}

It is pointless to read a book and learn \LaTeX.  Rather, you should
sit down and start writing. Use the books, the web, your peers and
your advisors as references when you get stuck.  \LaTeX\ can do almost
everything, and it does many things well.  
It is a free, portable standard
and the tool of choice for almost all mathematical journals and the
vast majority of technical journals.

\section{Important things that are not here}

Of course, there are many things are that are not in this document.
\begin{itemize}

\item Figures, tables and floats.

\item Bibliographical citations.  \LaTeX\ has a nice way of storing
  and using bibliographical databases called bib\TeX.

\item Equations covering more than one line.  Systems of equations
  grouped together.  Underbraces.  Overbraces.  Conditional
  statements.

\item Special environments like subsections, appendices, theorems and
  definitions.

\end{itemize}

\section{A final word}

\LaTeX\ makes everything look brilliant, even if it is total crud.  You
should not try to demonstrate your prowess by using every \LaTeX\ 
feature that you know.   Rather, you should be thoughtful about how
you write and then use \LaTeX \ to produce your writing.  In other
words, let
your ideas guide the tool, not the other way around.


\end{document}